\documentclass[11pt,a4paper]{amsart}
\pdfinfoomitdate=1
\pdftrailerid{}
\pdfsuppressptexinfo=15
\usepackage[T1]{fontenc}
\usepackage{lmodern}
\usepackage[margin=27mm]{geometry}
\usepackage{microtype}
\usepackage{amsmath,amssymb,mathtools}
\usepackage{booktabs}
\usepackage{xcolor}
\usepackage[numbers,sort&compress]{natbib}
\usepackage[colorlinks=true,linkcolor=blue!55!black,citecolor=blue!55!black,urlcolor=blue!55!black]{hyperref}
\usepackage[nameinlink,noabbrev]{cleveref}
\raggedbottom

\newtheorem{theorem}{Theorem}[section]
\newtheorem{proposition}[theorem]{Proposition}
\newtheorem{corollary}[theorem]{Corollary}
\newtheorem{lemma}[theorem]{Lemma}
\theoremstyle{definition}
\newtheorem{definition}[theorem]{Definition}
\theoremstyle{remark}
\newtheorem{remark}[theorem]{Remark}

\newcommand{\N}{\mathbb N}
\newcommand{\Z}{\mathbb Z}
\newcommand{\cO}{\mathcal O}
\newcommand{\supp}{\operatorname{supp}}
\newcommand{\End}{\operatorname{End}}
\newcommand{\Aut}{\operatorname{Aut}}
\newcommand{\Orb}{\operatorname{Orb}}
\newcommand{\Per}{\operatorname{Per}}
\newcommand{\htop}{h_{\mathrm{top}}}
\newcommand{\zz}{\mathbf 0}
\newcommand{\zero}{0^{\Z}}
\newcommand{\pq}{P_q}

\title[Digit-weight automatic towers]{Digit-Weight Automatic Towers:\\
Cantor--Bendixson Layers, Nilpotent Cellular Automata, and Conjugacy Rigidity}
\author{Anonymous}
\date{Internal preprint, 27 August 2026}
\subjclass[2020]{37B10, 68R15, 11B85, 54G12}
\keywords{automatic sequence, countable subshift, Cantor--Bendixson rank, cellular automaton, endomorphism monoid, conjugacy rigidity}

\begin{document}

\begin{abstract}
Fix integers $q\geq2$ and $d\geq1$, and mark the nonnegative integers whose
base-$q$ expansion has exactly $d$ digits equal to $1$ and every other digit
equal to $0$.  We determine the two-sided orbit closure of this indicator
sequence.  It is a countable transitive subshift consisting of one fixed
point and exactly $d+1$ free shift orbits, indexed by the number of retained
digits.  Its Cantor--Bendixson derivatives delete these orbits one at a time,
and its height is $d+2$.  Centered limits define continuous lowering cellular
automata $D_r$.  They generate, together with the shift, the full
endomorphism monoid:
\[
 \End(X_{q,d})=\{\zz\}\cup
 \{\sigma^aD_r:a\in\Z,\ 0\leq r\leq d\},
\]
with truncated-addition composition and an absorbing zero.  Consequently
$\Aut(X_{q,d})\cong\Z$.  We also prove arithmetic rigidity:
$X_{q,d}$ and $X_{p,e}$ are conjugate exactly when $(q,d)=(p,e)$.  The only
invariant probability is the point mass at the zero configuration, and the
entropy is zero.  The exponential-support case $d=1$ is treated as an owned
base case; the residual contribution is the exact package for $d\geq2$.
\end{abstract}

\maketitle

\section{Introduction}

For $q\geq2$ and $d\geq1$, let
\[
 S_{q,d}=\left\{\sum_{i\in A}q^i:
 A\subseteq\N_0,\ |A|=d\right\}.
\]
Thus $S_{q,d}$ consists of integers whose base-$q$ expansion has exactly
$d$ unit digits and no other nonzero digits.  Its indicator is $q$-automatic:
a finite automaton rejects digits outside $\{0,1\}$ and counts unit digits up
to $d$.  This is a direct instance of the classical automatic-sequence
framework initiated by Cobham \cite{Cobham1972}.  Sparse automatic supports
and their polylogarithmic growth have a substantial general theory; a recent
graph-theoretic account and precise ownership map are given by Wessel
\cite{Wessel2025}.

We study a different invariant: the complete two-sided orbit closure of the
indicator, including its topology and every continuous shift-commuting
self-map.  The answer is a finite tower of countably infinite orbits.  A
bounded difference of two fixed-cardinality sums of $q$-powers forces their
escaping exponent sets to agree.  This arithmetic rigidity classifies every
centered limit, including the binary case where an informal ``no carry''
argument would be insufficient.  The same classification makes the natural
layer-lowering maps continuous.

The endomorphism calculation is then forced by transitivity.  An endomorphism
is determined by the image of the generating point.  Every possible image is
realized by a shift followed by a lowering map, and the lowering maps form a
nilpotent chain.  The return scales from the top layer to the next layer are
the powers of $q$; requiring these scales in both directions under a
conjugacy recovers the base itself, including when the two bases are powers
of a common integer.

The case $d=1$ is not claimed as a new system.  Gheorghiciuc treats the exact
exponential occurrence function $G(n)=q^{n-1}$ and its word complexity
\cite{Gheorghiciuc2007}.  Salo and T\"orm\"a explicitly use the binary
powers-of-two orbit closure, with a slightly different initial-index
convention, as a one-dimensional gadget \cite{SaloTorma2012}.  General
automatic points in automatic systems and under factor maps are now
characterized through quasi-fixed points \cite{Krawczyk2026}.  We retain
$d=1$ in all theorem statements because it is the first level of the same
proof, but our novelty claims concern $d\geq2$.

The precise residual package is:
\begin{enumerate}
\item the exact orbit decomposition and Cantor--Bendixson tower;
\item continuous lowering cellular automata and the full endomorphism
monoid, not merely the automorphism group;
\item conjugacy rigidity in both parameters $(q,d)$;
\item the complete periodic-point and invariant-measure description.
\end{enumerate}
A bounded primary-source search through 27 August 2026, covering the exact
defining equation, fixed-digit-weight automatic sets, powers-of-two orbit
closures, countable automatic systems, and endomorphism calculations, found
no source stating this combined $d\geq2$ package.  This is a bounded search
report, not a claim of absolute priority.

\section{Setup and the main theorem}

Write $\N_0=\{0,1,2,\ldots\}$.  For a finite set $A\subseteq\N_0$, put
\[
 \pq(A)=\sum_{i\in A}q^i.
\]
The map $\pq$ is injective because its values have base-$q$ digits in
$\{0,1\}$.  Set $S_{q,0}=\{0\}$ and define $x_{q,s}\in\{0,1\}^{\Z}$ by
\[
 x_{q,s}(n)=\begin{cases}
 1,&n\in S_{q,s},\\
 0,&n\notin S_{q,s}.
 \end{cases}
\]
In particular, $x_{q,0}$ is the single marker at the origin.  We use the
left shift
\[
 (\sigma^a x)(n)=x(n+a),
 \qquad \supp(\sigma^a x)=\supp(x)-a.
\]
Let
\[
 X_{q,d}=\overline{\{\sigma^a x_{q,d}:a\in\Z\}},
 \qquad
 \cO_s=\{\sigma^a x_{q,s}:a\in\Z\}.
\]
The symbol $\zero$ denotes the all-zero configuration.  We reserve $\zz$ for
the constant map with value $\zero$.

By an endomorphism we mean a continuous map $F:X_{q,d}\to X_{q,d}$ satisfying
$F\sigma=\sigma F$; surjectivity is not part of the definition.  Hedlund's
theorem \cite{Hedlund1969} identifies these maps with sliding block codes on
the subshift.

\begin{theorem}[main theorem]\label{thm:main}
For $q\geq2$ and $d\geq1$, the following hold.
\begin{enumerate}
\item The orbit decomposition is the disjoint union
\begin{equation}\label{eq:orbit-decomposition}
 X_{q,d}=\{\zero\}\sqcup\bigsqcup_{s=0}^{d}\cO_s.
\end{equation}
\item For $0\leq j\leq d$,
\begin{equation}\label{eq:derivatives}
 X_{q,d}^{(j)}=\{\zero\}\cup
 \bigcup_{s=0}^{d-j}\cO_s,
\end{equation}
while $X_{q,d}^{(d+1)}=\{\zero\}$ and
$X_{q,d}^{(d+2)}=\varnothing$.  Hence the Cantor--Bendixson height is
$d+2$.
\item There are continuous lowering maps $D_r$, $0\leq r\leq d$, such that
\begin{equation}\label{eq:lowering-definition}
 D_r(\zero)=\zero,
 \qquad
 D_r(\sigma^a x_{q,s})=
 \begin{cases}
 \sigma^a x_{q,s-r},&s\geq r,\\
 \zero,&s<r.
 \end{cases}
\end{equation}
They satisfy $D_0=\mathrm{id}$ and
\begin{equation}\label{eq:lowering-composition}
 D_rD_t=
 \begin{cases}
 D_{r+t},&r+t\leq d,\\
 \zz,&r+t>d.
 \end{cases}
\end{equation}
\item The full endomorphism monoid is
\begin{equation}\label{eq:end-monoid}
 \End(X_{q,d})=\{\zz\}\cup
 \{\sigma^aD_r:a\in\Z,\ 0\leq r\leq d\}.
\end{equation}
Its multiplication is
\begin{equation}\label{eq:end-multiplication}
 (\sigma^aD_r)(\sigma^bD_t)=
 \begin{cases}
 \sigma^{a+b}D_{r+t},&r+t\leq d,\\
 \zz,&r+t>d.
 \end{cases}
\end{equation}
Consequently $\Aut(X_{q,d})=\langle\sigma\rangle\cong\Z$.
\item If $p\geq2$ and $e\geq1$, then
\begin{equation}\label{eq:conjugacy-rigidity}
 (X_{q,d},\sigma)\cong(X_{p,e},\sigma)
 \quad\Longleftrightarrow\quad (q,d)=(p,e).
\end{equation}
\item The only periodic point is $\zero$, the unique invariant Borel
probability is $\delta_{\zero}$, and $\htop(X_{q,d})=0$.
\end{enumerate}
\end{theorem}

The restriction $d,e\geq1$ in the rigidity statement is necessary.  At
$d=0$, the orbit closure of the single marker is independent of $q$.

\section{Bounded differences of digital sums}

The following lemma is the arithmetic engine of every limit argument below.

\begin{lemma}[bounded-difference rigidity]\label{lem:bounded-difference}
Fix $q\geq2$ and $M\geq0$.  Let $(A_j)$ and $(B_j)$ be sequences of finite
subsets of $\N_0$ with $|A_j|,|B_j|\leq M$.  Suppose that
\[
 |\pq(A_j)-\pq(B_j)|\leq C
\]
for a constant $C$ independent of $j$.  After passing to a subsequence,
there are fixed finite sets $A,B\subseteq\N_0$ and finite sets $E_j$ such
that
\begin{equation}\label{eq:common-escape}
 A_j=A\sqcup E_j,
 \qquad B_j=B\sqcup E_j,
 \qquad \min E_j\longrightarrow\infty,
\end{equation}
where the last condition is vacuous when $E_j$ is empty.
\end{lemma}

\begin{proof}
The indicator functions of subsets of $\N_0$ lie in the compact product
space $\{0,1\}^{\N_0}$.  A diagonal extraction makes membership of every
fixed exponent eventually constant in both sequences.  Let $A$ and $B$ be
the respective pointwise limits.  They are finite and have cardinality at
most $M$: otherwise $M+1$ distinct exponents would eventually belong to one
of the approximating sets.

After discarding finitely many indices, write
\[
 A_j=A\sqcup A'_j,
 \qquad B_j=B\sqcup B'_j,
\]
where the least exponent in $A'_j\cup B'_j$ tends to infinity whenever that
union is nonempty.  Remove the common part
$E_j=A'_j\cap B'_j$, and put
\[
 U_j=A'_j\setminus E_j,
 \qquad V_j=B'_j\setminus E_j.
\]
The sets $U_j$ and $V_j$ are disjoint.  Moreover,
\[
 |\pq(U_j)-\pq(V_j)|
 \leq C+\pq(A)+\pq(B)=:C'.
\]
If $U_j\cup V_j$ is nonempty, let $m_j$ be its least exponent.  Then
$\pq(U_j)-\pq(V_j)$ is divisible by $q^{m_j}$.  For large $j$, one has
$q^{m_j}>C'$, so this difference must be zero.  Injectivity of $\pq$ gives
$U_j=V_j$; disjointness then forces both sets to be empty.  Thus
$A'_j=B'_j=E_j$ for all sufficiently large $j$, proving
\eqref{eq:common-escape}.
\end{proof}

The proof uses divisibility only after every stable low exponent has been
removed.  It therefore covers $q=2$ without assuming that subtraction is
carry-free.

\section{Centered limits and orbit decomposition}

\begin{proposition}[centered-limit normal form]\label{prop:centered-limits}
Fix $0\leq s\leq d$.  Suppose
\[
 \sigma^{n_j}x_{q,s}\longrightarrow y\neq\zero.
\]
After passing to a subsequence, there exist $t\in\{0,\ldots,s\}$, an integer
$a$, and sets $E_j\subseteq\N_0$ such that
\begin{equation}\label{eq:center-normal-form}
 |E_j|=s-t,
 \qquad \min E_j\longrightarrow\infty,
 \qquad n_j=a+\pq(E_j),
\end{equation}
and
\begin{equation}\label{eq:center-limit}
 y=\sigma^a x_{q,t}.
\end{equation}
Conversely, every sequence satisfying \eqref{eq:center-normal-form}
converges as in \eqref{eq:center-limit}.
\end{proposition}

\begin{proof}
Choose $k_0\in\Z$ with $y(k_0)=1$.  Coordinatewise convergence in the
discrete alphabet implies that, for all sufficiently large $j$,
\[
 n_j+k_0=\pq(A_j)
\]
for a unique $s$-element set $A_j$.  Pass to a subsequence on which the
membership of every fixed exponent stabilizes.  There is then a fixed finite
set $B$ and sets $E_j$ whose least exponents tend to infinity such that
\[
 A_j=B\sqcup E_j.
\]
After one further extraction, $|E_j|$ is constant.  Put
$t=|B|$ and $a=\pq(B)-k_0$.  This gives
\eqref{eq:center-normal-form}.

We identify the support of $y$.  If $a+k=\pq(C)$ for a $t$-element set
$C$, then $C\cap E_j=\varnothing$ for large $j$, and
\[
 n_j+k=\pq(E_j\sqcup C)\in S_{q,s}.
\]
Hence $y(k)=1$.  Conversely, suppose $y(k)=1$.  For large $j$ there are
$s$-element sets $C_j$ satisfying
\[
 \pq(C_j)=n_j+k,
 \qquad
 \pq(C_j)-\pq(A_j)=k-k_0.
\]
Apply \cref{lem:bounded-difference} to $(C_j)$ and $(A_j)$.  The stable
part of $(A_j)$ is $B$, so the lemma gives a fixed $t$-element set $C$ and,
after relabelling the common escaping part, $C_j=C\sqcup E_j$.  Therefore
\[
 \pq(C)-\pq(B)=k-k_0,
 \qquad \pq(C)=a+k.
\]
We have proved
\[
 y(k)=1\quad\Longleftrightarrow\quad a+k\in S_{q,t},
\]
which is exactly $y=\sigma^a x_{q,t}$.

For the converse, assume \eqref{eq:center-normal-form}.  If
$a+k\in S_{q,t}$, adjoining its unique $t$ exponents to $E_j$ proves that
$n_j+k\in S_{q,s}$ for large $j$.  If $n_j+k\in S_{q,s}$ along an infinite
subsequence, compare the representing $s$-element sets with $E_j$ by
\cref{lem:bounded-difference}; it follows that $a+k\in S_{q,t}$.  Thus every
coordinate converges to that of $\sigma^a x_{q,t}$.
\end{proof}

\begin{proof}[Proof of \cref{thm:main}(1)]
Every nonzero limit of shifts of $x_{q,d}$ belongs to one of the displayed
orbits by \cref{prop:centered-limits}.  The zero configuration also occurs:
$\sigma^{-j}x_{q,d}\to\zero$ as $j\to\infty$, since the support of
$x_{q,d}$ is contained in $\N_0$.  Conversely, for $0\leq t\leq d$, choose
$(d-t)$-element sets $E_j$ whose least exponent tends to infinity.
\Cref{prop:centered-limits} gives
\[
 \sigma^{\pq(E_j)}x_{q,d}\longrightarrow x_{q,t},
\]
and shifting this convergence produces every point of $\cO_t$.

It remains to verify disjointness.  Every nonzero point in $\cO_t$ has
counting function
\[
 \#\bigl(\supp(x)\cap[-N,N]\bigr)=\Theta((\log N)^t)
\]
for $t\geq1$, while points of $\cO_0$ have singleton support.  Translation
does not change the exponent of this polylogarithmic growth, so different
layers cannot meet.  Indeed, the upper bound follows because every exponent
in a representation of a point of $S_{q,t}\cap[0,N]$ is at most
$\lfloor\log_q N\rfloor$, while the lower bound follows by choosing all $t$
exponents below $\lfloor\log_q(N/t)\rfloor$; injectivity of $\pq$ makes these
choices distinct.  Within one layer, a nonzero period would make a
nonempty support that is bounded below invariant under a nonzero
translation, which is impossible.  Thus every $\cO_t$ is a free orbit.
\end{proof}

\section{The Cantor--Bendixson tower}

Let
\[
 Y_s=\{\zero\}\cup\bigcup_{t=0}^{s}\cO_t,
 \qquad 0\leq s\leq d.
\]

\begin{lemma}[one orbit removed at a time]\label{lem:one-orbit}
For $s\geq1$, the derived set of $Y_s$ is $Y_{s-1}$, while the derived set
of $Y_0$ is $\{\zero\}$ and the derived set of $\{\zero\}$ is empty.
\end{lemma}

\begin{proof}
Every point $\sigma^a x_{q,t}$ with $t<s$ is approached by points of
$\cO_s$: choose $(s-t)$-element escaping sets $E_j$ and use
\[
 \sigma^{a+\pq(E_j)}x_{q,s}\longrightarrow\sigma^a x_{q,t}.
\]
The zero configuration is an accumulation point of every $Y_s$: for
$s=0$, translate the single marker to infinity; for $s\geq1$, use
$\sigma^{-j}x_{q,s}\to\zero$.

We show that every point of $\cO_s$ is isolated in $Y_s$.  Otherwise a
sequence of distinct points of $Y_s$ would converge to a point of $\cO_s$.
After taking a subsequence, all points lie in one layer $\cO_t$, $t\leq s$.
If their shift parameters are unbounded, \cref{prop:centered-limits} says
that every nonzero limit belongs to a layer of index strictly below $t$.
It therefore cannot lie in $\cO_s$.  If the parameters are bounded, a
subsequence is constant; freeness of the orbit then contradicts distinctness.
This proves $(Y_s)'=Y_{s-1}$.

The same argument for $Y_0$ says that all single-marker points are isolated
and their shifts accumulate only at $\zero$.  Finally, the singleton
$\{\zero\}$ has empty derived set.
\end{proof}

\begin{proof}[Proof of \cref{thm:main}(2)]
Since $X_{q,d}=Y_d$, iterating \cref{lem:one-orbit} gives
\eqref{eq:derivatives}, followed by $\{\zero\}$ and then the empty set.  In
the convention where the height is the least ordinal with empty derivative,
the height is $d+2$.  Equivalently, it is the least stable derivative index.
The point rank of every member of $\cO_s$ is $d-s$, and the point rank of
$\zero$ is $d+1$.
\end{proof}

For example, when $d=1$ the derivative chain is
\[
 X_{q,1}\supset \cO_0\cup\{\zero\}
 \supset\{\zero\}\supset\varnothing.
\]
This makes the otherwise easy off-by-one error visible.

\section{Lowering maps and their continuity}

The orbit decomposition makes \eqref{eq:lowering-definition} a well-defined
set map.  Shift commutation is immediate.  Its continuity is the nonlocal
point of the construction.

\begin{proposition}[continuity of lowering]\label{prop:lowering-continuity}
For each $0\leq r\leq d$, the map $D_r$ in
\eqref{eq:lowering-definition} is continuous.  Hence it is a cellular
automaton on $X_{q,d}$.
\end{proposition}

\begin{proof}
Because $X_{q,d}$ is compact metrizable, it suffices to prove sequential
continuity.  Let $y_j\to y$.  There are only finitely many layers, so every
subsequence has a further subsequence which either consists entirely of
$\zero$ or lies in a fixed nonzero layer.  The former case is immediate.  In
the latter case, write
\[
 y_j=\sigma^{n_j}x_{q,s}.
\]
We prove that its $D_r$-images have a further subsequence converging to
$D_r(y)$.  This subsequence criterion implies convergence of the full image
sequence.

First suppose $y=\sigma^a x_{q,t}$ is nonzero.  By
\cref{prop:centered-limits}, after a further extraction the centers have
the form
\[
 n_j=a+\pq(E_j),
 \qquad |E_j|=s-t,
 \qquad \min E_j\to\infty.
\]
If $s<r$, then $t\leq s<r$ and both $D_r(y_j)$ and $D_r(y)$ are zero.  Assume
$s\geq r$.  If $t\geq r$, then
\[
 D_r(y_j)=\sigma^{n_j}x_{q,s-r}
 \longrightarrow \sigma^a x_{q,t-r}=D_r(y)
\]
by the converse half of \cref{prop:centered-limits}.

If $t<r$, we claim that $D_r(y_j)\to\zero$.  Were there a nonzero sublimit,
\cref{prop:centered-limits} would express the same centers as
\[
 n_j=b+\pq(F_j),
 \qquad |F_j|\leq s-r,
 \qquad \min F_j\to\infty.
\]
The difference $\pq(E_j)-\pq(F_j)=b-a$ is bounded.  Applying
\cref{lem:bounded-difference} to the two escaping families forces
$|E_j|=|F_j|$ eventually.  This is impossible because
\[
 |E_j|=s-t>s-r\geq|F_j|.
\]
Thus the only sublimit is $\zero=D_r(y)$.

It remains to consider $y=\zero$.  If $D_r(y_j)$ had a nonzero sublimit,
then necessarily $s\geq r$ (the case $s<r$ is identically zero), and
the same normal form would give
\[
 n_j=b+\pq(F_j),
 \qquad |F_j|=h\leq s-r
\]
after a further extraction.  The sets $F_j$ escape to infinity.
Applying the converse part of \cref{prop:centered-limits} to the original
layer $s$ would then give the nonzero limit
$\sigma^b x_{q,s-h}$, where $s-h\geq r$, contradicting
$y_j\to\zero$.  Therefore $D_r(y_j)\to\zero$.

The map is continuous and commutes with $\sigma$.  Hedlund's theorem now
supplies a finite-radius local rule, so $D_r$ is a cellular automaton.
\end{proof}

\begin{proof}[Proof of \cref{thm:main}(3)]
Only the composition law remains.  On a layer $\cO_s$, applying $D_t$ and
then $D_r$ lowers the index by $r+t$ when $s\geq r+t$ and gives $\zero$
otherwise.  This is $D_{r+t}$ when $r+t\leq d$.  If $r+t>d$, every layer is
sent to $\zero$, so the composite is $\zz$.
\end{proof}

In particular, $D_1^{d+1}=\zz$.  The law is truncation to an absorbing zero,
not saturation at the bottom nonzero layer.

\section{The full endomorphism monoid}

\begin{proof}[Proof of \cref{thm:main}(4)]
The orbit of $x_{q,d}$ is dense in $X_{q,d}$.  Consequently two continuous
shift-commuting maps that agree on $x_{q,d}$ agree on its orbit and hence on
all of $X_{q,d}$.

Let $F\in\End(X_{q,d})$.  If $F(x_{q,d})=\zero$, then
$F(\sigma^a x_{q,d})=\zero$ for every $a$, so continuity and density give
$F=\zz$.  Otherwise \eqref{eq:orbit-decomposition} gives unique
$a\in\Z$ and $0\leq s\leq d$ such that
\[
 F(x_{q,d})=\sigma^a x_{q,s}.
\]
The endomorphism $\sigma^aD_{d-s}$ has the same value at $x_{q,d}$, hence
$F=\sigma^aD_{d-s}$.  This proves that \eqref{eq:end-monoid} contains every
endomorphism; \cref{prop:lowering-continuity} proves the reverse inclusion.

The displayed maps are distinct.  Different lowering indices send
$x_{q,d}$ to different layers.  At a fixed lowering index, different shift
exponents give different points because every nonzero orbit is free.  The
constant map is separate.  Since the lowering maps commute with the shift,
\eqref{eq:end-multiplication} follows from
\eqref{eq:lowering-composition}.

For $r>0$, the image of $\sigma^aD_r$ misses the top $r$ layers; $\zz$ is
also not invertible.  Every $\sigma^aD_0=\sigma^a$ is invertible.  The units
are therefore precisely the shift powers, proving
$\Aut(X_{q,d})\cong\Z$.
\end{proof}

\begin{remark}[surjective convention]
Some authors reserve ``endomorphism'' for onto shift-commuting maps.  Under
that convention the preceding proof shows that the surjective endomorphism
monoid equals the automorphism group $\langle\sigma\rangle$.  Throughout this
paper, $\End$ uses the broader cellular-automaton convention.
\end{remark}

\section{Conjugacy rigidity}

We first isolate the arithmetic consequence of a one-digit degeneration.

\begin{lemma}[one-digit return scales]\label{lem:one-digit-scales}
Let $p\geq2$, $s\geq1$, and $u_j\to+\infty$.  If
\[
 \sigma^{u_j}x_{p,s}\longrightarrow\sigma^b x_{p,s-1},
\]
then for all sufficiently large $j$ there is an exponent $m_j\in\N_0$ such
that
\begin{equation}\label{eq:one-digit-centers}
 u_j=b+p^{m_j}.
\end{equation}
\end{lemma}

\begin{proof}
Every subsequence has, by \cref{prop:centered-limits}, a further subsequence
of the form $u_j=a+p^{m_j}$ whose limit is $\sigma^a x_{p,s-1}$.  Freeness of
$\cO_{s-1}$ forces $a=b$.  If \eqref{eq:one-digit-centers} failed along
infinitely many indices, applying this argument to precisely those indices
would give a contradiction.
\end{proof}

\begin{lemma}[two-sided exponential matching]\label{lem:exponential-matching}
Let $p,q\geq2$.  Suppose that for all sufficiently large $N$ there are
integers $m_N$ and a fixed integer $K$ such that
\begin{equation}\label{eq:exponential-difference}
 q^N-p^{m_N}=K.
\end{equation}
Then $q=p^u$ for an integer $u\geq1$, and in fact $K=0$ and $m_N=uN$
eventually.  If the analogous conclusion also holds with $p$ and $q$
interchanged, then $p=q$.
\end{lemma}

\begin{proof}
The right side of $p^{m_N}=q^N-K$ is eventually strictly increasing, so
$m_N$ is eventually strictly increasing.  Dividing consecutive equations
gives
\[
 p^{m_{N+1}-m_N}
 =\frac{q^{N+1}-K}{q^N-K}\longrightarrow q.
\]
The left side lies in the discrete set $\{p,p^2,\ldots\}$.  It is therefore
eventually constant and equal to $q$, so $q=p^u$ for some $u\geq1$.  Next,
\[
 p^{m_N-uN}=\frac{p^{m_N}}{q^N}\longrightarrow1.
\]
This is a sequence in the discrete set $p^{\Z}$, hence it eventually equals
$1$.  Thus $m_N=uN$ and \eqref{eq:exponential-difference} gives $K=0$.

If also $p=q^v$ for an integer $v\geq1$, then
$q=q^{uv}$.  Since $q>1$, one has $uv=1$, so $u=v=1$ and $p=q$.
\end{proof}

\begin{proof}[Proof of \cref{thm:main}(5)]
The forward implication is immediate when $(q,d)=(p,e)$.  Conversely, let
$H:X_{q,d}\to X_{p,e}$ be a conjugacy.  Cantor--Bendixson height is a
topological invariant, so $d+2=e+2$ and hence $d=e$.

The isolated points of $X_{q,d}$ are exactly $\cO_d$, while the points of
rank one are exactly $\cO_{d-1}$.  Thus there are integers $b,c$ with
\[
 H(x_{q,d})=\sigma^c x_{p,d},
 \qquad
 H(x_{q,d-1})=\sigma^b x_{p,d-1}.
\]
For every $N$,
\[
 \sigma^{q^N}x_{q,d}\longrightarrow x_{q,d-1}.
\]
Applying $H$ and using shift commutation gives
\[
 \sigma^{q^N+c}x_{p,d}\longrightarrow\sigma^b x_{p,d-1}.
\]
By \cref{lem:one-digit-scales},
\[
 q^N+c=b+p^{m_N}
\]
for all sufficiently large $N$.  Hence
$q^N-p^{m_N}=b-c$ and \cref{lem:exponential-matching} gives $q=p^u$.
Applying the same argument to $H^{-1}$ gives $p=q^v$.  The last part of
\cref{lem:exponential-matching} yields $p=q$.
\end{proof}

The inverse conjugacy is essential.  For example, every power of $4$ is a
power of $2$, but the odd powers of $2$ have no matching return scale in base
$4$.

\section{Periodic points, invariant measures, and entropy}

\begin{proof}[Proof of \cref{thm:main}(6)]
The point $\zero$ is fixed.  Every other point has a nonempty support bounded
below.  If such a point had nonzero period $k$, its support would be invariant
under translation by $k$ and would consequently be unbounded in both
directions.  This contradiction proves
\[
 \Per(X_{q,d})=\{\zero\}.
\]

The space $X_{q,d}$ is countable by \eqref{eq:orbit-decomposition}.  Let
$\mu$ be an invariant Borel probability.  All atoms on a fixed infinite
orbit have the same mass.  That mass must be zero, since otherwise the orbit
would have infinite total mass.  There are only finitely many such orbits,
so their union has measure zero.  Therefore
$\mu(\{\zero\})=1$ and $\mu=\delta_{\zero}$.

The measure-theoretic entropy of $\delta_{\zero}$ is zero.  The variational
principle for continuous maps of compact metric spaces now gives
$\htop(X_{q,d})=0$.
\end{proof}

The system is therefore topologically transitive but uniquely ergodic on a
fixed point; its transitive point is not recurrent.  The nontrivial structure
is topological and functorial rather than measure-theoretic.

\section{Ownership boundary and reproducible checks}

The sequence $x_{q,1}$ has support $1,q,q^2,\ldots$.  Gheorghiciuc's
exponential occurrence function is exactly this one-sided word and yields
linear factor complexity \cite{Gheorghiciuc2007}.  In the binary case,
Salo and T\"orm\"a call the orbit closure with markers at $2^n$, $n>0$, the
binary powers-of-two shift \cite{SaloTorma2012}.  Their starting convention
differs by one initial marker from ours.  These sources prevent any claim
that the $d=1$ seed, exponential sparsity, or zero entropy is new.

Cobham's theory owns automaticity \cite{Cobham1972}; Wessel and the sources
therein own the general sparse-support dichotomy and automatic rank
\cite{Wessel2025}; Krawczyk owns the general characterization of automatic
points in an automatic system and its factors \cite{Krawczyk2026}.  None of
those results is restated here as a contribution.  Our residual object is
the $d\geq2$ tower, and the residual claims are the exact centered-limit
normal form, Cantor--Bendixson filtration, full nilpotent endomorphism monoid,
and two-parameter conjugacy rigidity.

The accompanying script \texttt{code/verify\_digit\_weight.py} checks finite
shadows of the fragile statements.  It tests many centered windows with
nonconsecutive escaping exponents and nonzero phases, the absorbing
truncated-addition law, a divisibility form of
\cref{lem:bounded-difference}, and bidirectional return-scale separation.
These computations are deterministic controls, not substitutes for the
all-scale proofs.

\section{Conclusion}

Fixed unit-digit weight converts an elementary regular language into a
countable transitive subshift with a rigid finite topological tower.  The
same escaping-digit mechanism controls the orbit closure, every derivative,
every cellular automaton, and every conjugacy.  The resulting endomorphism
monoid is the product of integer shifts with a finite nilpotent lowering
chain, completed by an absorbing zero.  The construction gives such a tower
at every finite height at least three, while the base $q$ remains visible in
the one-digit return scales.

Two limitations are deliberate.  First, the paper treats only unit digits;
allowing arbitrary nonzero digit values introduces genuine carry relations.
Second, the ownership statement is bounded to the primary-source search
described above and does not assert absolute priority.  Both extensions are
natural directions for further work.

\bibliographystyle{plainnat}
\bibliography{references}

\end{document}
